% problem-000145 4 化学反应与结构综合分析
\section*{4. 化学反应与结构综合分析}
（图：）

（图：）
图 4-1

\subsection*{4-1}
为了调节荧光蛋⽩的吸收光谱，可以对其氨基酸序列进⾏改进。⼀种改进⽅法是将酪氨酸突变为组氨酸，其吸收峰发⽣蓝移。画出对应的发⾊团结构简式（不要求⽴体构型）。

\subsection*{4-2}
另⼀种改进⽅式是引⼊⼀个新的酪氨酸，其侧链残基恰好位于原有的酪氨酸上⽅，结构如图4-2所⽰。该发⾊团吸收峰与绿⾊荧光蛋⽩相⽐发⽣了红移，简述原因。

（图：）
图 4-2

\subsection*{4-3}
⼀种红⾊荧光蛋⽩其发⾊团胆绿素经由⾎红素氧化⽽产⽣，同时释放亚铁离⼦。反应过程如图4-3所⽰。给出亚铁离⼦离去的主要原因。

（图：）
图 4-3

\subsection*{4-4}
绿⾊荧光蛋⽩在pH=7.40下总电荷数为−7.88，其氨基酸序列中侧链残基带电荷的所有氨基酸种类以及他们的侧链 $) K _ { a }$ 如下图所⽰（可忽略⽔的电离的影响，和形成多肽对侧链 $\mathsf { p } K _ { a }$ 的影响）。已知绿⾊荧光蛋⽩的氨基酸序列中赖氨酸数⽬为20个，天冬氨酸数⽬为18个，⾕氨酸数⽬为16个，精氨酸数⽬ $\left( n _ { \mathrm { R } } \right)$ ⼩于10个，组氨酸的数⽬ $( n _ { \mathrm { H } } )$ ⼤于5个，组氨酸和半胱氨酸的数⽬之和 $( n _ { \mathrm { H } } + n _ { \mathrm { C } } )$ 为12个。计算 $n _ { \mathrm { R } } .$ 、 $n _ { \mathrm { H } } ^ { \mathrm { } } .$ 和� 值。

（图：）
图 4-4
